% !TeX root = test-main.tex

\begin{center}
{\Large\textbf{第4回 確認テスト}}
\end{center}
\vspace{0.15em}
\begin{flushright}
氏名：\testnameblank
\end{flushright}
\vspace{0.25em}
\noindent 以下の問題は、配布テキストに沿って解答すること。
\vspace{0.35em}

\noindent\textbf{1. 比較級、最上級}(5 $\times$ 10点)

\noindent\textbf{【問1】} 比較(比較級・最上級・原級) できるのは何か。簡単に答えなさい。
\testblank[形容詞または副詞]{0.92\linewidth}
\testqsep

\noindent\textbf{【問2】} 次の単語の比較級及び最上級を答えなさい。

good, bad, little
\testblank[better - best, worse - worst, less - least]{0.92\linewidth}
\testqsep

\noindent\textbf{【問3】} 比較級の基本構文とその意味を簡単に答えなさい。
\testblank[A + 比較級 + than + B \qquad 「AよりもB」]{0.92\linewidth}
\testqsep

\noindent\textbf{【問4】} 最上級の基本構文において、of と in の使い分けを簡単に答えなさい。
\testblank[of: 範囲を表す語句 \qquad in: 場所を表す語句]{0.92\linewidth}
\testqsep

\noindent\textbf{【問5】} 次の日本語を英文にしなさい。

「この問題はあの問題より簡単だ。」
\testblank[This problem is easier than that one.]{0.92\linewidth}
\testqsep

\noindent\textbf{2. 原級}(3 $\times$ 10点)

\noindent\textbf{【問6】} 原級のつくり方を簡単に答えなさい。
\testblank[変化させず as と as で形容詞・副詞をはさむ]{0.92\linewidth}
\testqsep

\noindent\textbf{【問7】} 次の英文を和訳しなさい。

Tom is not as tall as Mike.
\testblank[トムはマイクほど背が高くない]{0.92\linewidth}
\testqsep

\noindent\textbf{3. 比較}(3 $\times$ 10点)

\noindent\textbf{【問8】} 比較の否定文において、注意しなければならないことを簡単に答えなさい。
\testblank[使われている動詞に着目して、否定文の構造を決定すること]{0.92\linewidth}
\testqsep

\noindent\textbf{【問9】} 次の英文を意味を変えないように書き換えなさい。

No other person is as smart as Riku.
\testblank[Riku is the smartest of all.]{0.92\linewidth}
\testqsep

\noindent\textbf{【問10】} 次の英文に対する答えを主語と動詞を含む英文1文で答えなさい。

Which city has the largest population in Hokkaido?
\testblank[Sapporo is the largest city in Hokkaido.]{0.92\linewidth}
\testqsep

\begin{flushright}
\textbf{得点}\testscoreblank ／ 100点
\end{flushright}
